\chapter*{Tóm tắt}
\addcontentsline{toc}{chapter}{Tóm tắt}

Sự phát triển nhanh chóng của các hệ thống mạng viễn thông hiện đại đã đặt ra nhiều thách thức lớn trong quá trình vận hành và xử lý sự cố. Các phương pháp chẩn đoán truyền thống thường dựa vào luật thủ công hoặc các mô hình AI tĩnh, gặp nhiều khó khăn trong việc thích ứng với các lỗi mới hoặc sự thay đổi của cấu trúc mạng. Nhằm giải quyết vấn đề này, đề tài tập trung nghiên cứu, xây dựng và đánh giá hệ thống **Self-Evolving AI Agents** hỗ trợ tự động xử lý sự cố mạng dựa trên nền tảng **NIKA** (A Network Arena for Benchmarking AI Agents on Network Troubleshooting).

Nền tảng NIKA cung cấp môi trường mô phỏng mạng thực tế dựa trên Kathará, hỗ trợ giả lập nhiều sự cố phức tạp (như lỗi cấu hình định tuyến BGP/OSPF, lỗi thiết bị mạng, SDN controller và P4 data-plane) cùng cơ chế đánh giá hiệu năng chẩn đoán tự động. Nghiên cứu này hiện thực hóa hai mô-đun cốt lõi giúp AI Agent có khả năng tự tiến hóa:
\begin{enumerate}
    \item \textbf{Mô-đun Bộ Nhớ Quy Trình Kỹ Năng (Skill-Pro Memory Module):} Lưu trữ các kỹ năng chẩn đoán dạng \texttt{ProceduralSkill} trong file JSON, kết hợp cơ chế truy xuất nhận biết thuộc tính ngữ cảnh lấy cảm hứng từ MemInsight và tiến hóa kỹ năng ngoại tuyến thông qua Semantic Gradient và PPO Gate phi tham số (lấy cảm hứng từ LightMem và A-Mem). Quá trình tiến hóa được kiểm soát bằng ngưỡng phần thưởng kết hợp (detection, localization, RCA) nhằm chỉ học hỏi từ các episode chẩn đoán có chất lượng cao.
    \item \textbf{Mô-đun Tiến Hóa Tài Liệu Công Cụ (DRAFT Tool Evolution Module):} Hỗ trợ AI Agent tự tinh chỉnh mô tả (documentation) của các công cụ MCP nguyên thủy tại test-time, lấy cảm hứng từ DRAFT. Vòng lặp Explorer--Analyzer--Rewriter liên tục trích xuất trial thực tế, xác định khoảng trống hiểu biết (ComprehensionGap) và viết lại tài liệu chi tiết hơn. Tài liệu tự động bị đóng băng khi hội tụ, đảm bảo hiệu quả tài nguyên dài hạn.
\end{enumerate}

Hiệu quả của hệ thống được kiểm chứng thông qua các thực nghiệm chẩn đoán chi tiết trên tập con 30 sự cố chọn lọc từ benchmark NIKA. Kết quả cho thấy agent được trang bị mô-đun tiến hóa đạt độ chính xác chẩn đoán vượt trội so với baseline thông thường (điểm F1 chẩn đoán tăng từ 0.0111 lên 0.1611, cải thiện đáng kể khả năng chẩn đoán trong các sự cố mất cấu hình định tuyến và IP), đồng thời giảm số lượng bước gọi công cụ thừa và tối ưu hóa chi phí token sử dụng.

\vspace{10pt}
\noindent \textit{\textbf{Từ khóa}: Self-Evolving AI Agents, xử lý sự cố mạng, NIKA, bộ nhớ quy trình tiến hóa, tiến hóa công cụ, Kathará.}
